\documentclass[11pt,a4paper]{article}
\usepackage[T1]{fontenc}
\usepackage[utf8]{inputenc}
\usepackage{lmodern,microtype}
\usepackage[margin=20mm,headheight=14pt]{geometry}
\usepackage{amsmath,amssymb,booktabs,tabularx,array,enumitem,xcolor,titlesec,fancyhdr}
\definecolor{navy}{RGB}{25,51,77}
\titleformat{\section}{\large\bfseries\color{navy}}{}{0pt}{}
\titlespacing*{\section}{0pt}{8pt}{4pt}
\setlength{\parindent}{0pt}
\setlength{\parskip}{5pt}
\setlist[enumerate]{leftmargin=17pt,itemsep=3pt,topsep=3pt}
\pagestyle{fancy}
\fancyhf{}
\fancyhead[L]{\small Master's thesis project}
\fancyhead[R]{\small Deep learning and anomaly detection}
\fancyfoot[C]{\small\thepage}
\renewcommand{\headrulewidth}{0.3pt}
\newcommand{\arxiv}[1]{arXiv:#1}
\newcommand{\lead}[1]{\textbf{#1}\ }
\begin{document}
\fontsize{10.5}{12.6}\selectfont
{\Large\bfseries\color{navy} Deep learning for anomaly detection\\in graviton-inspired collider searches\par}

\section{1.\ Physics motivation and project objective}

Quantising weak gravitational fields gives a massless spin-2 graviton
whose ordinary gravitational coupling is too weak for realistic direct
production at the LHC \cite{gravity}. Randall--Sundrum (RS) models
instead predict additional massive spin-2 excitations that can decay
into two photons \cite{rs,atlas}. This project will quantify
sensitivity to a small set of such $G_{\mathrm{KK}}\to\gamma\gamma$
benchmarks. Standard Model diphoton production, jets misidentified as
photons, detector effects and simulation mismodelling must all be
considered. An unusual event population alone does not establish a
gravitational origin.

The research question is: \emph{How do reconstruction, generative
modelling and limited supervision affect sensitivity to rare signals
and generalisation to unseen signal parameters?}  The core will
compare a deep autoencoder, a variational autoencoder and Deep SAD,
with simple kinematic selections as a reference. At most one further
architecture or generative method will be developed, depending on data
access and progress. Hjorth-Jensen's textbook provides the theoretical
and computational preparation \cite{book}.

\section{2.\ First phase: CERN educational open data}
ATLAS educational resources at the CERN Open Data Portal provide reconstructed collision events, simulated samples and guided notebooks \cite{data}. The student will use a notebook and dataset from the same release, beginning with small samples before scaling up. The 2025 release offers updated educational tools; the legacy 2020 diphoton collection provides a documented $10\,\mathrm{fb}^{-1}$ example at $13$ TeV.
\begin{enumerate}
\item \textbf{Understand the data.} Read ROOT files with Python/Uproot or the supplied framework. Inspect photons, leptons and jets; check units, transverse momenta, angles, triggers, event weights and preselection.
\item \textbf{Reconstruct known physics.} Reproduce the $Z$ mass using a dilepton sample, then work through the Higgs-to-diphoton notebook. Calculate invariant masses, compare normalised distributions and study how selections change signal and background.
\item \textbf{Prepare a learning benchmark.} Use simulated Higgs and background events to test reconstruction scores, then reveal a small subset of training labels for semi-supervised learning. Examine the ATLAS data-quality example to understand instrumental anomalies.
\end{enumerate}

This phase produces a reproducible notebook with selection counts,
distributions, uncertainties and an initial benchmark. The legacy
release is educational and not suitable for scientific publications but should give a good introduction to the data to come for your thesis. The
main study requires suitable research data and compatible
signal/background simulations, or an explicit simulation-only setup
with a common detector model.

\section{3.\ Deep-learning methods}
\lead{Autoencoders: reconstruction as the baseline (textbook Ch.~12 and covered by FYS5429).}
For standardised event features $x$, an encoder compresses the event to $z=f_\theta(x)$ and a decoder reconstructs $\hat{x}=g_\psi(z)$. Training minimises the average reconstruction loss on background-dominated events; a simple anomaly score is
\[
 a_{\mathrm{AE}}(x)=\|x-g_\psi(f_\theta(x))\|^2 .
\]
A bottleneck or denoising objective discourages copying the input. Start with fully connected networks for high-level photon/event variables \cite{ae}. Large errors can flag unfamiliar structure, but neither anomalous signals nor detector artefacts are guaranteed to be poorly reconstructed. This is the label-free baseline.

\lead{Variational autoencoders: a generative extension (Ch.~15 of the textbook and covered by FYS5429).}
A VAE learns an approximate latent posterior $q_\phi(z|x)$ and a probabilistic decoder $p_\theta(x|z)$ with prior $p(z)$, commonly Gaussian \cite{vae,collidervae}. Its per-event negative evidence lower bound is
\[
 \mathcal{L}_{\mathrm{VAE}}(x)=
 -\mathbb{E}_{q_\phi(z|x)}[\log p_\theta(x|z)]
 +D_{\mathrm{KL}}\!\left(q_\phi(z|x)\,\|\,p(z)\right).
\]
The terms balance reconstruction and latent regularity. Sampling the prior and decoding generates events. Reconstruction and negative-bound scores require validation; the bound is not an exact likelihood or anomaly probability. Standard VAEs use no labels; an added labelled classification loss enables semi-supervision.

\newpage
\lead{Deep SAD: using a few labels to shape the representation.}
Deep Semi-supervised Anomaly Detection (Deep SAD) \cite{sad} fine-tunes an encoder, typically pretrained with an autoencoder. Its score is the squared distance to a fixed latent centre,
$a_{\mathrm{SAD}}(x)=\|f_\theta(x)-c\|^2$.
Predominantly normal unlabelled events and labelled background are pulled towards $c$; labelled simulated signals are pushed away. Writing $a=a_{\mathrm{SAD}}$, a representative regularised objective is
\[
 \mathcal{L}_{\mathrm{SAD}}
 =\langle a(x)\rangle_U+\lambda_B\langle a(x)\rangle_B
 +\lambda_A\left\langle\frac{1}{a(x)+\delta}\right\rangle_A
 +\lambda_R\|\theta\|^2 .
\]
Here $U,B,A$ denote unlabelled, labelled background and labelled anomalous samples; brackets denote sample averages, $\lambda$ values control their weights, and $\delta>0$ stabilises the inverse. The implementation will follow the published safeguards against collapse, including a fixed nonzero centre, bias-free layers and suitable unbounded activations. Varying label budgets and contamination tests whether supervision helps; withheld signal parameters test generalisation. This changes the encoder's training objective rather than introducing a separate architecture.

\lead{Connections and optional extensions.}
A transformer specifies how features interact; a generative objective specifies what a model learns. They can therefore be combined. The options below will be assessed before selecting at most one extension.

{\small
\renewcommand{\arraystretch}{1.13}
\begin{tabularx}{\linewidth}{@{}>{\raggedright\arraybackslash\bfseries}p{30mm}X@{}}
\toprule
Method & \textbf{Role in this project and conditions for use}\\
\midrule
Transformers (Ch.~13 and FYS5429) &
Self-attention learns relations between particle/object tokens \cite{transformer}. A set encoder with permutation-invariant readout could replace a dense encoder or learn masked-object prediction. This is most motivated by richer object/constituent data; an advantage is not assumed for a few diphoton variables.\\
Normalising flows (Ch.~16 and FYS5429) &
Invertible neural transformations provide tractable densities and sampling. A conditional background model $p_B(x|m_{\gamma\gamma})$, learned in mass sidebands, could generate templates for a data-versus-background classifier, following the CATHODE strategy \cite{flow}. Interpolation and background closure require validation.\\
Diffusion models (Ch.~16 and FYS5429) &
A network learns to reverse progressive noising \cite{diffusion}. Possible extensions are denoising-based scores or conditional background generation, including in a VAE latent space. Noise levels, sampling cost and rare-tail fidelity require study; generation alone does not define a calibrated anomaly score.\\
Reinforcement learning &
A possible later project concerns sequential trigger or data-acquisition decisions under a resource budget. This requires a simulator, actions and a meaningful reward. The present fixed-dataset anomaly search has no natural sequential control task, so RL is outside the core scope.\\
\bottomrule
\end{tabularx}
}

\section{4.\ Comparison and statistical validation}

All core methods will use the same event partitions and comparable
tuning budgets. Preprocessing is fitted on training data; validation
selects hyperparameters and thresholds; the test set remains
untouched. Labels hidden from AE/VAE training are revealed selectively
for Deep SAD. The study varies signal contamination, labelled sample
size and withheld masses/couplings. Independent physically simulated
signal injections begin during model development. Generated
backgrounds complement these tests; they do not independently
establish sensitivity to new signals.

Report signal efficiency at fixed background efficiency, ROC and
precision--recall curves, uncertainty and computational
cost. Precision depends on the assumed signal fraction; confusion
matrices require truth-labelled samples. Reserve diphoton mass for
inference and test score--mass correlations: excluding mass from
inputs alone does not prevent artificial peaks. Check detector/model
variations, simulation-to-data differences and background predictions
in control regions.

After fixing the selection, use a mass-spectrum likelihood with
nuisance parameters for background and detector uncertainties
\cite{stats}. Calibrate with background-only pseudoexperiments,
repeating any data-dependent fitting, and account for scans when
quoting global significance. Anomaly scores select events; statistical
evidence concerns an ensemble excess. Deliver expected sensitivity or
justified limits, including a careful account of failures and
generalisation.

\newpage
\section{5.\ Milestones}

Educational exercises begin during M1--M3 and are consolidated in
M4. Signal injections start at M5 and continue throughout; the
optional extension follows successful core comparisons.

{\small
\renewcommand{\arraystretch}{1.13}
\begin{tabularx}{\linewidth}{@{}>{\bfseries}p{8mm}X@{}}
\toprule
& \textbf{Milestone and completion criterion}\\
\midrule
M1 & Study collider observables, backgrounds, statistics and neural networks (textbook Ch.~8); run an introductory ATLAS notebook.\\
M2 & Select the spin-2 final state, benchmark masses and couplings; compile physics and deep-learning references.\\
M3 & Study and implement small AE/VAE examples (Chs.~12, 15); derive their losses and understand the Deep SAD objective.\\
M4 & Complete the $Z$, diphoton and data-quality exercises; document variables, weights and selections; choose the main data/simulation route.\\
M5 & Fix event partitions; establish kinematic and deep-AE baselines with background-only and injected-signal tests.\\
M6 & Train VAE and Deep SAD models; compare scores, label budgets and contamination using common validation criteria.\\
M7 & Check stability, mass sculpting and background closure; assess one optional transformer, flow or diffusion extension if feasible.\\
M8 & Quantify efficiencies, uncertainties, computational costs and expected statistical sensitivity, including scan effects.\\
M9 & Test withheld signal parameters and detector/modelling variations; assess generalisation and interpret the RS benchmark sensitivity.\\
M10 & Complete injection/recovery studies; deliver reproducible code, notebooks, configurations, figures and thesis conclusions.\\
\bottomrule
\end{tabularx}
}

\lead{Preparation and expected outcome.}

The textbook reading progresses from neural networks and autoencoders
to variational inference, followed by attention or diffusion/flows for
the chosen extension \cite{book}. The thesis will determine which
representation and objective improve sensitivity under realistic
limitations. No improvement, limited transfer to unseen signals, or a
quantified sensitivity limit are meaningful outcomes.

\begingroup
\fontsize{8.8}{10.2}\selectfont
\setlength{\parskip}{0pt}
\renewcommand{\refname}{6.\ Literature and data resources}
\begin{thebibliography}{13}
\setlength{\itemsep}{2.5pt}
\bibitem{gravity}
J.~F.~Donoghue, \emph{General relativity as an effective field theory: The leading quantum corrections},
Phys.\ Rev.\ D \textbf{50}, 3874 (1994), \arxiv{gr-qc/9405057}.
\bibitem{rs}
L.~Randall and R.~Sundrum, \emph{A Large Mass Hierarchy from a Small Extra Dimension},
Phys.\ Rev.\ Lett.\ \textbf{83}, 3370 (1999), \arxiv{hep-ph/9905221}.
\bibitem{atlas}
ATLAS Collaboration, \emph{Search for resonances decaying into photon pairs in 139 fb$^{-1}$ of $pp$ collisions at $\sqrt{s}=13$ TeV with the ATLAS detector},
Phys.\ Lett.\ B \textbf{822}, 136651 (2021), \arxiv{2102.13405}.
\bibitem{book}
M.~Hjorth-Jensen, \emph{Machine Learning and the Physical Sciences}, University of Oslo,
MachineLearningUiO online textbook, Chs.~8, 12, 13, 15 and 16.
\bibitem{data}
ATLAS Collaboration, \emph{ATLAS Open Data}, educational and research resources; 2025 educational release; CERN Open Data records 15006 (legacy diphotons) and atlas-160012 (data quality).
\bibitem{ae}
M.~Farina, Y.~Nakai and D.~Shih, \emph{Searching for New Physics with Deep Autoencoders},
Phys.\ Rev.\ D \textbf{101}, 075021 (2020), \arxiv{1808.08992}.
\bibitem{vae}
D.~P.~Kingma and M.~Welling, \emph{Auto-Encoding Variational Bayes},
ICLR (2014), \arxiv{1312.6114}.
\bibitem{collidervae}
O.~Cerri et al., \emph{Variational Autoencoders for New Physics Mining at the Large Hadron Collider},
JHEP \textbf{05}, 036 (2019), \arxiv{1811.10276}.
\bibitem{sad}
L.~Ruff et al., \emph{Deep Semi-Supervised Anomaly Detection}, ICLR (2020),
\arxiv{1906.02694}.
\bibitem{transformer}
H.~Qu, C.~Li and S.~Qian, \emph{Particle Transformer for Jet Tagging},
PMLR \textbf{162}, 18281 (2022), \arxiv{2202.03772}.
\bibitem{flow}
A.~Hallin et al., \emph{Classifying Anomalies THrough Outer Density Estimation (CATHODE)},
Phys.\ Rev.\ D \textbf{106}, 055006 (2022), \arxiv{2109.00546}.
\bibitem{diffusion}
J.~Ho, A.~Jain and P.~Abbeel, \emph{Denoising Diffusion Probabilistic Models},
NeurIPS \textbf{33} (2020), \arxiv{2006.11239}.
\bibitem{stats}
G.~Cowan et al., \emph{Asymptotic formulae for likelihood-based tests of new physics},
Eur.\ Phys.\ J.\ C \textbf{71}, 1554 (2011), \arxiv{1007.1727}.
\end{thebibliography}
\endgroup
\end{document}

